\begin{titlepage}

\begin{figure}[H]
  \begin{center}
    \includegraphics[width=3cm]{auth.pdf}
    \label{fig:cover_auth_logo}
  \end{center}
\end{figure}

\centering
\Large Αριστοτέλειο Πανεπιστήμιο Θεσσαλονίκης\\
\Large Πολυτεχνική Σχολή\\
\large Τμήμα Ηλεκτρολόγων Μηχανικών και Μηχανικών Υπολογιστών\\
\large Τομέας Ηλεκτρονικής \& Υπολογιστών

\vspace{\fill}

\LARGE Ψηφιακά Δίδυμα στην Εκπαίδευση: Ανάπτυξη Συστήματος \textlatin{IoT} για τη Βελτιστοποίηση της Ενεργειακής Απόδοσης και του Μαθησιακού Περιβάλλοντος

\vspace{\fill}

\Large Διπλωματική Εργασία\\
\Large του\\
\Large Αλέξανδρου Παπαδόπουλου

\vspace{\fill}
\raggedright

\begin{tabular}{ll}
\textbf{Επιβλέπων καθηγητής:} & Αλκιβιάδης Χατζόπουλος, Α.Π.Θ.\\
\end{tabular}

\centering
\vspace{\fill}
\today

\end{titlepage}
\setcounter{page}{0}

\chapter*{Περίληψη}
\addcontentsline{toc}{chapter}{Περίληψη}

Η παρούσα διπλωματική εργασία πραγματεύεται τον σχεδιασμό, την υλοποίηση και την αξιολόγηση μιας ολοκληρωμένης πλατφόρμας Ψηφιακού Διδύμου (\textlatin{Digital Twin}), ειδικά προσαρμοσμένης για εκπαιδευτικά περιβάλλοντα και Έξυπνες Πανεπιστημιουπόλεις (\textlatin{Smart Campuses}). Σε αντίθεση με τις παραδοσιακές, αντιδραστικές (\textlatin{reactive}) μεθόδους διαχείρισης κτιρίων ή τα απλά συστήματα παρακολούθησης (\textlatin{Digital Shadows}), η προτεινόμενη λύση εγκαθιδρύει μια αμφίδρομη, πραγματικού χρόνου επικοινωνία μεταξύ του φυσικού σχολικού χώρου και του ψηφιακού του αντιγράφου, επιτρέποντας τη λήψη αυτοματοποιημένων αποφάσεων.

Η ανάπτυξη του συστήματος ακολούθησε μια αυστηρή διεπιστημονική προσέγγιση, βασισμένη στην αρχιτεκτονική πέντε επιπέδων. Σε επίπεδο υλικού, εγκαταστάθηκε δίκτυο κόμβων Διαδικτύου των Πραγμάτων (\textlatin{IoT}) χαμηλού κόστους, αποτελούμενο από μικροελεγκτές \textlatin{ESP32}, έξυπνους μετρητές \textlatin{Shelly 1PM} και συσκευές \textlatin{Sonoff}. Σε επίπεδο λογισμικού, αναπτύχθηκε μια ισχυρή, τοπικά εκτελούμενη αρχιτεκτονική με χρήση \textlatin{MQTT, InfluxDB} και \textlatin{openHAB} ως σημασιολογικό ενδιάμεσο λογισμικό (\textlatin{edge middleware}) στο τοπικό επίπεδο, το οποίο τροφοδοτεί ένα \textlatin{Enterprise Cloud SaaS} επίπεδο (\textlatin{React 18 / Supabase}) που υλοποιεί την τελική τρισδιάστατη (\textlatin{3D}) διεπαφή χρήστη.

Η έρευνα δομείται πάνω σε τρεις κεντρικούς πυλώνες: (α) τη βελτιστοποίηση του μαθησιακού περιβάλλοντος (\textlatin{Environmental Quality}), μέσω της διαρκούς παρακολούθησης της ακουστικής άνεσης με επιχειρησιακό όριο παρέμβασης τα \textlatin{55 dBA}, (β) την ενεργειακή αποδοτικότητα (\textlatin{Energy Efficiency}), με συνεχή καταγραφή ισχύος, και (γ) τη λειτουργική διαχείριση (\textlatin{Facility Management}) μέσω απομακρυσμένου ελέγχου. Από την πιλοτική εφαρμογή στα «Εκπαιδευτήρια Πλάτων» (Μάρτιος -- Νοέμβριος 2025), επιτεύχθηκε και επαληθεύτηκε μέση μείωση της ενεργειακής κατανάλωσης κατά 19,7\%. Παράλληλα, η τεχνοοικονομική ανάλυση επιβεβαίωσε την υψηλή προστιθέμενη αξία της μετάβασης σε ένα σύγχρονο Έξυπνο Σχολείο, καταδεικνύοντας απόσβεση της επένδυσης σε 42,5 μήνες και 5ετές \textlatin{ROI} 41,02\%.

\noindent\textbf{Λέξεις-Κλειδιά:} Ψηφιακό Δίδυμο, Έξυπνο Σχολείο, Διαδίκτυο των Πραγμάτων (\textlatin{IoT}), Μαθησιακό Περιβάλλον, Ενεργειακή Αποδοτικότητα, Διαχείριση Εγκαταστάσεων.

\clearpage
\selectlanguage{english}
\chapter*{Abstract}
\addcontentsline{toc}{chapter}{Abstract}

This diploma thesis presents the design, implementation, and evaluation of a comprehensive \textlatin{Digital Twin} platform tailored specifically for educational environments and \textlatin{Smart Campuses}. In contrast to traditional reactive building-management methods or simple monitoring systems (\textlatin{Digital Shadows}), the proposed solution establishes a bidirectional, real-time communication channel between the physical school environment and its digital counterpart, enabling automated decision-making.

The system was developed through a rigorous interdisciplinary approach based on a five-layer architecture. At the hardware level, a low-cost \textlatin{Internet of Things} (\textlatin{IoT}) network was deployed, consisting of ESP32 microcontrollers, Shelly 1PM smart meters, and Sonoff devices. At the software level, a locally executed architecture was implemented using the MQTT protocol, the InfluxDB time-series database, and the open-source openHAB middleware operating as edge semantic middleware, which feeds an Enterprise Cloud SaaS layer (React 18 / Supabase) that delivers the final three-dimensional (3D) user interface.

The research is structured around three core pillars: (\textlatin{a}) optimization of the learning environment (\textlatin{Environmental Quality}) through continuous acoustic-comfort monitoring based on a 55 dBA intervention threshold; (\textlatin{b}) \textlatin{Energy Efficiency} through continuous power tracking; and (\textlatin{c}) \textlatin{Facility Management} through centralized monitoring and remote actuation. Following the pilot implementation at Platon Schools (March -- November 2025), an average energy-consumption reduction of 19.7\% was achieved and verified. In parallel, the techno-economic analysis confirmed the added value and financial viability of the proposed approach, demonstrating a 42.5-month payback period and a five-year ROI of 41.02\%.

\noindent\textbf{Keywords:} \textlatin{Digital Twin}, \textlatin{Smart Campus}, \textlatin{Internet of Things} (\textlatin{IoT}), Learning Environment, \textlatin{Energy Efficiency}, \textlatin{Facility Management}.
\selectlanguage{greek}
\clearpage


\chapter*{Ευχαριστίες}
\addcontentsline{toc}{chapter}{Ευχαριστίες}
Η παρούσα διπλωματική εργασία αποτελεί το επιστέγασμα μιας μακράς, απαιτητικής, αλλά και εξαιρετικά δημιουργικής
προσπάθειας. Η ολοκλήρωσή της δεν θα ήταν εφικτή χωρίς την πολύτιμη συμβολή, καθοδήγηση και υποστήριξη ανθρώπων
που στάθηκαν δίπλα μου και πίστεψαν στο όραμα αυτής της πλατφόρμας.

Αρχικά, θα ήθελα να εκφράσω τις θερμότερες και ειλικρινείς ευχαριστίες μου στον επιβλέποντα καθηγητή μου,
κ. Άλκη Χατζόπουλο, για την εμπιστοσύνη που μου έδειξε, την πολύτιμη ακαδημαϊκή του καθοδήγηση και τη συνεχή
υποστήριξη καθ' όλη τη διάρκεια εκπόνησης της εργασίας.

Ένα ξεχωριστό και μεγάλο ευχαριστώ οφείλω στον κ. Χρήστο Κάλλα. Η συμβολή του υπήρξε καθοριστική, καθώς
αποτέλεσε την αφετηρία και την τεχνική βάση ολόκληρης της έρευνας. Η βοήθειά του στην εγκατάσταση του υλικού
εξοπλισμού (\textlatin{hardware}) και των αισθητήρων \textlatin{IoT}, καθώς και η πολύτιμη καθοδήγησή του στο στήσιμο και την
παραμετροποίηση της πλατφόρμας \textlatin{openHAB}, αποτέλεσαν το θεμέλιο πάνω στο οποίο αναπτύχθηκε το παρόν Ψηφιακό
Δίδυμο.

Τέλος, θα ήθελα να ευχαριστήσω θερμά τη διοίκηση, τους εκπαιδευτικούς και το προσωπικό των «Εκπαιδευτηρίων
Πλάτων» στην Κατερίνη. Η προθυμία τους να αποδεχθούν το σχολείο τους ως πεδίο μελέτης περίπτωσης
(\textlatin{case study}) και να φιλοξενήσουν την πιλοτική εφαρμογή του συστήματος, μου έδωσε τη μοναδική ευκαιρία να
εφαρμόσω τη θεωρία στην πράξη σε ένα πραγματικό, ζωντανό εκπαιδευτικό περιβάλλον.

\clearpage
